\section*{Заключение}
\addcontentsline{toc}{section}{Заключение}

В рамках данной работы было исследовано математическое основание для объемного рендеринга, используемого в задаче визуализации трехмерных воксельных данных, среди которых можно выделить несколько основных идей: формула альфа блендинга (\ref{eq:alpha_blending1}, \ref{eq:alpha_blending2}) и отпечатка трехмерных объектов на двумерной плоскости, если лучи выходят из нее перпендикулярно (\ref{eq:footprint}).

Эти идеи использует современная технология 3D Gaussian Splatting, которая стала основным объектом исследования. Она предлагает развить идеи прямого объемного рендеринга и сделать его дифференцируемым, представив сцену в виде поля излучения: для позиции в $\mathcal R^3$ и направления взгляда из $S^2$ выдающее непрозрачность и цвет. Для этого получен способ рассчета отпечатка гауссианы в случае, когда лучи исходят из некоторого оптического центра камеры: с некоторым приближением он считается гауссианой с ковариацией (первые 2 строки и столбца из формулы \ref{eq:sigma_footprint}). Также было рассмотрено какие параметры стоит оптимизировать и какие применить дополнительные методики обучения (прунинг и уплотнение).

К стандартному конвейеру 3DGS, состоящий в формировании из облака точек и камер COLMAP в набор гауссиан, было разработано множество скриптов, которые:
%
\begin{itemize}
	\item Генерируют из воксельных данных КТ (и в текущей работе RGB данных согласованных с КТ) заданное количество рендеров трассировки путей методом Монте-Карло (его реализует движок Cycles от Blender);
	\item Далее запускается стандартный алгоритм 3DGS, это возможно благодаря тому, что предыдущий этап записывает свои результаты сразу в формат COLMAP;
	\item Полученные гауссианы сегментируются по методу ближайшего соседа с размеченной воксельной сеткой, и эти данные добавляются к стандартному формату как дополнительное поле;
	\item Визуализацией сегментированных гауссиан занимается WebGL, где для оптимизации из них вырезаны сферические гармоники, отвечающие за зависимый от направления взгляда цвет.
\end{itemize}
%
Для демонстрации данным пайплайном была построена сегментированная костная система человека. В результате которого получены 700 реалистичных изображений и сегментированные гауссианы.

Полученный пайплайн несовершенен, далее описаны перспективы развития:
%
\begin{itemize}
	\item Добавить к этапу рендеринга поверхности, так как уравнение переноса излучения не учитывает волновой природы света. Для этого потребуется выделять поверхности и приписывать им коэффициенты приломления;
	\item Заменить RGB воксельную сетку на функции переноса, так как цветные изображения это феноменальные данные, присущие только используемому датасету, на практике имеются лишь данные КТ;
	\item Добавить автоматическое сегментирование органов;
	\item Заменить стандартный 3DGS на более продвинутый аналоги, например \cite{jiang2023gaussianshader3dgaussiansplatting}. Аналогично первому замечанию, гауссианы сами по себе строятся на теории объемного рендеринга, что делает их оригинальную реализацию ограниченной. Например, зеркальные отражения в референсных изображениях 3DGS реализует как просто дублирование сцены за поверхностью зеркала. Тем самым у итоговой модели теряются полезные свойства, например нормали к поверхности.
\end{itemize}